\documentclass[11pt]{article}
\usepackage[utf8]{inputenc}
\usepackage[T1]{fontenc}
\usepackage{amsmath,amssymb}
\usepackage{graphicx}
\usepackage{booktabs}
\usepackage{hyperref}
\usepackage[margin=1in]{geometry}
\usepackage{natbib}
\bibliographystyle{unsrtnat}

\title{Thiolutin Has Complex Effects \textit{In Vivo} but Is a Direct Inhibitor of RNA Polymerase II \textit{In Vitro}}
\author{Authors}
\date{}

\begin{document}
\maketitle

\begin{abstract}
Thiolutin is a natural product transcription inhibitor with an unresolved mode of action. Thiolutin and the related dithiolopyrrolone holomycin chelate Zn$^{2+}$ and previous studies have concluded that RNA Polymerase II (Pol II) inhibition \textit{in vivo} is indirect. Here, we present chemicogenetic and biochemical approaches to investigate thiolutin's mode of action in \textit{Saccharomyces cerevisiae}. We identify mutants that alter sensitivity to thiolutin. We provide genetic evidence that thiolutin causes oxidation of thioredoxins \textit{in vivo} and that thiolutin both induces oxidative stress and interacts functionally with multiple metals including Mn$^{2+}$ and Cu$^{2+}$, and not just Zn$^{2+}$. Finally, we show direct inhibition of RNA polymerase II (Pol II) transcription initiation by thiolutin \textit{in vitro} in support of classical studies that thiolutin can directly inhibit transcription \textit{in vitro}. Inhibition requires both Mn$^{2+}$ and appropriate reduction of thiolutin as excess DTT abrogates its effects. Pause-prone, defective elongation can be observed \textit{in vitro} if inhibition is bypassed. Thiolutin effects on Pol II occupancy \textit{in vivo} are widespread but major effects are consistent with prior observations for Tor pathway inhibition and stress induction, suggesting that thiolutin use \textit{in vivo} should be restricted to studies on its modes of action and not as an experimental tool.
\end{abstract}

\section{Introduction}

Thiolutin is a historically used transcription inhibitor with potentially multiple modes of action. Thiolutin is one of the microbially produced dithiolopyrrolone compounds that feature an intra-molecular and redox-sensitive disulfide bond \citep{Tipper1973,Oliva2001}. Holomycin is another notable compound in this class. These compounds are thought to be pro-drugs whose anti-microbial action requires reduction of the intramolecular disulfide. Thiolutin has been shown to inhibit bacterial and eukaryotic transcription \textit{in vivo} and has been used to study mRNA stability in multiple species \citep{Jimenez1973,Parker1991,Grigull2004,Pelechano2010}.

Recent studies have produced conflicting results regarding thiolutin's mechanism of action. While classical \textit{in vitro} experiments demonstrated direct Pol II inhibition by thiolutin \citep{Tipper1973}, more recent work has concluded that Pol II inhibition is indirect, possibly mediated through Zn$^{2+}$ chelation affecting the proteasome or other cellular targets \citep{Lauber2021,Chan2017}. The discovery that reduced thiolutin and holomycin chelate Zn$^{2+}$ has provided one mechanism for their cellular effects \citep{Chan2017,Lauber2021}. However, whether Zn$^{2+}$ chelation fully accounts for transcription inhibition remains unresolved.

Here, we use a combination of chemicogenetic screening and biochemistry to dissect thiolutin's mode of action in \textit{S. cerevisiae}. Our genetic screens reveal roles for multidrug resistance, oxidative stress response, and metal homeostasis pathways in thiolutin sensitivity. Critically, we demonstrate that thiolutin, when activated by reduction and Mn$^{2+}$, directly inhibits Pol II transcription initiation \textit{in vitro}, reconciling conflicting reports and revealing why previous studies lacking Mn$^{2+}$ in their buffers failed to observe direct inhibition.

\section{Results}

\subsection{Thiolutin or Reduced Thiolutin Alone Fail to Inhibit Purified Pol II \textit{In Vitro}}

Recent publications suggest that thiolutin does not directly inhibit Pol II, in contrast to classic \textit{in vitro} transcription experiments \citep{Tipper1973}. Given that thiolutin and holomycin activities could be sensitive to reduction of the disulfide bond by DTT \citep{Tipper1973,Chan2017,Lauber2021} (Figure~\ref{fig:structure}A), we removed DTT from the \textit{in vitro} transcription reaction and performed thiolutin treatment with or without molar equivalents of DTT (Figure~\ref{fig:structure}B). Under the buffer conditions employed, we failed to observe thiolutin inhibition of Pol II, with DTT having no effect. We therefore set out to screen for potential co-factors.

\subsection{Three Independent Genetic Screens for Thiolutin Resistant and Sensitive Mutants}

To gain insights into factors controlling the cellular response to thiolutin, we performed three independent genetic screens for thiolutin-resistant and -sensitive mutants (Figure~\ref{fig:screens}A). First, we performed a forward genetic screen by UV mutagenesis, identifying causal mutations by bulk segregant analysis and whole-genome sequencing. Second, we screened a yeast Variomics library for individual thiolutin-resistant candidates \citep{Pan2004}. Third, we performed high-throughput screens using both pooled Variomics and deletion libraries using barcode sequencing (Bar-seq).

Our data exhibited excellent reproducibility and revealed valuable common and distinct features between Variomics (putative gain-of-function and loss-of-function mutants combined) and deletion (loss-of-function mutants) libraries. We identified contributions of multidrug resistance (MDR), oxidative stress response (OSR), multiple divalent metal homeostasis/trafficking pathways, and the proteasome to thiolutin resistance.

\subsection{Multidrug Resistance and Oxidative Stress Pathways}

From both manual Variomics and forward genetic screens, we isolated multiple mutants in YAP1 and YRR1, encoding transcription factors implicated in OSR and/or MDR. In addition, we isolated mutants in the MDR efflux pump gene SNQ2 from the manual Variomics screen. Isolated yap1 mutations were clustered in C-terminal regions, where a variety of mutants are known to cause gain-of-function. MDR mutants from the manual Variomics screen were dominant or functioned through increased dosage.

In Bar-seq screens, Variomics mutants in several MDR genes conferred thiolutin resistance whereas deletion of the same genes either conferred sensitivity or had no strong effect, consistent with these Variomics mutants being different from gene-loss-of-function. We conclude that the MDR transcription factor Pdr1 is the major transcription factor promoting efflux of thiolutin from the cell.

\subsection{Thiolutin Causes Oxidation of Thioredoxins}

Thioredoxins are small anti-oxidant proteins that primarily function in reducing specific cysteines in client proteins. In yeast, there are two thioredoxin reductases (Trr1 and Trr2) and three thioredoxins (Trx1, Trx2, Trx3). Mutants in TRR1 were reproducibly isolated from our manual Variomics screen, and trr1$\Delta$ conferred hypersensitivity to thiolutin, consistent with Trr1 activity antagonizing thiolutin.

Critically, trx1$\Delta$ trx2$\Delta$ conferred thiolutin resistance and suppressed thiolutin hypersensitivity of trr1$\Delta$, consistent with the hypothesis that trr1$\Delta$ hypersensitivity is due to accumulation of oxidized thioredoxins. This reveals a critical requirement of Trx1 and Trx2 for thiolutin activity \textit{in vivo}.

Thiolutin induced Yap1::EGFP nuclear localization, consistent with induction of oxidative stress, and depleted total glutathione within one hour of treatment. Thiolutin could be reduced by DTT or TCEP, but not equimolar glutathione, and reduced thiolutin could be re-oxidized, suggesting thiolutin acts as a redox-cycler. However, thiolutin did not cause DNA damage as measured by mutation frequency at CAN1, and treatment was reversible.

\subsection{Thiolutin Alters Divalent Metal Homeostasis}

We observed distinct thiolutin resistance or sensitivity linked to diverse Zn$^{2+}$ trafficking genes, including ZAP1 and various Zn$^{2+}$ transporters. zap1$\Delta$ conferred hypersensitivity to both thiolutin and holomycin, consistent with Zn$^{2+}$ chelation being a conserved mechanism among dithiolopyrrolones. Zn$^{2+}$ supplementation partially suppressed thiolutin inhibition but did not fully rescue it, suggesting additional mechanisms beyond Zn$^{2+}$ chelation.

In addition to Zn$^{2+}$ homeostasis mutants, we observed genes involved in Mn$^{2+}$ trafficking linked to thiolutin resistance. pmr1$\Delta$, which increases intracellular Mn$^{2+}$, conferred thiolutin hypersensitivity. Mn$^{2+}$ supplementation dose-dependently exacerbated thiolutin sensitivity in wild-type yeast (Figure~\ref{fig:metals}C), supporting synergy between Mn$^{2+}$ and thiolutin \textit{in vivo}.

We tested whether reduced thiolutin chelates Mn$^{2+}$ and Cu$^{2+}$, while confirming Zn$^{2+}$ chelation. Consistent with recent reports, Zn$^{2+}$ caused a distinct UV shift of reduced thiolutin. We found that Mn$^{2+}$ and Cu$^{2+}$ can also interact with reduced thiolutin, while Mg$^{2+}$ did not bind (Figure~\ref{fig:metals}E). $^1$H-NMR experiments confirmed interactions between the divalent metals and reduced thiolutin.

\subsection{Thiolutin, When Activated by DTT and Mn$^{2+}$, Directly Inhibits Pol II \textit{In Vitro}}

Remarkably, Mn$^{2+}$ addition to reduced thiolutin potently inhibited purified Pol II activity (Figure~\ref{fig:inhibition}A), in sharp contrast to the activation of Pol II by Mn$^{2+}$ alone. We note that the original Tipper observations used transcription buffer containing 1.6 mM Mn$^{2+}$ \citep{Tipper1973}, whereas most subsequent studies did not include Mn$^{2+}$ in their buffers \citep{Lauber2021,Chan2017}.

Template binding protects Pol II from Thio/Mn$^{2+}$ inhibition (Figure~\ref{fig:inhibition}B), suggesting inhibition at an early step in transcription consistent with transcription initiation inhibition. On a transcription bubble template where initiation is bypassed, Thio/Mn$^{2+}$ altered but did not block Pol II elongation, inducing a highly pause-prone mode (Figure~\ref{fig:inhibition}E). Pre-assembled elongation complexes were also resistant to Thio/Mn$^{2+}$.

Thio/Mn$^{2+}$ inhibition required appropriate reduction: excess DTT (12 mM) abolished inhibition (Figure~\ref{fig:inhibition}C--D). UV spectroscopy revealed dynamic changes in species after Mn$^{2+}$ addition to reduced thiolutin, with a 380 nm peak species forming within 2 minutes and subsequently converting to a 300 nm species (Figure~\ref{fig:species}A). Freshly prepared Thio/Mn$^{2+}$ strongly inhibited Pol II, but the inhibitory activity was completely lost after 20 minutes (Figure~\ref{fig:species}B), consistent with the 380 nm species contributing to Pol II inhibition yet being unstable. However, immediately treated Pol II remained inhibited for at least 50 minutes, suggesting stable Pol II modification.

\subsection{Thiolutin Reduces Pol II Occupancy Genome-Wide}

ChIP-qPCR showed specific decrease of Pol II occupancy at the 5' end of YLR454W within the first 2 minutes of thiolutin treatment, consistent with fast transcription initiation inhibition \textit{in vivo} (Figure~\ref{fig:chipseq}A). The relatively slower loss of Pol II from the template compared to other initiation inhibition conditions was consistent with additional, possibly direct effects on elongation.

Pol II ChIP-seq with spike-in normalization showed broad decrease in Pol II occupancy genome-wide (Figure~\ref{fig:chipseq}C). We observed a continuum of effects where ribosomal protein (RP) genes were affected most strongly, then ribosome biogenesis (RiBi) genes, then other genes (Figure~\ref{fig:chipseq}D--E). These classes of effect likely involve Tor signaling inhibition and stress response activation, consistent with thiolutin having complex \textit{in vivo} effects beyond direct Pol II inhibition.

\section{Discussion}

Our study reconciles conflicting reports about thiolutin's mode of action by identifying Mn$^{2+}$ as a critical co-factor for direct Pol II inhibition. The original demonstration of thiolutin-mediated transcription inhibition by Tipper used buffers containing Mn$^{2+}$ \citep{Tipper1973}, but subsequent studies that concluded thiolutin does not directly inhibit Pol II omitted this metal from their reaction conditions. This explains the discrepancy in the literature.

The Thio/Mn$^{2+}$ complex inhibits transcription initiation, as shown by the protection conferred by template pre-binding and the bypass of inhibition on elongation templates. When inhibition is bypassed, Thio/Mn$^{2+}$-treated Pol II shows pause-prone elongation, suggesting that the complex may also interact with Pol II during elongation but that this interaction is insufficient to completely block transcription.

Our genetic screens reveal that thiolutin has pleiotropic effects \textit{in vivo}, involving MDR, oxidative stress, thioredoxin oxidation, and perturbation of multiple divalent metal homeostasis pathways. The widespread but differential effects on Pol II occupancy genome-wide, with RP genes being most affected, are consistent with Tor pathway inhibition and general stress responses. These complex \textit{in vivo} effects suggest that thiolutin should not be used as a simple transcription shutoff tool for mRNA stability studies, but rather should be studied for its own modes of action.

The transient nature of the inhibitory Thio/Mn$^{2+}$ species (active for less than 20 minutes in solution but stably modifying Pol II) suggests a reactive intermediate that may form a covalent or coordination complex with Pol II. The suppression by excess DTT is consistent with involvement of sulfenylation or disulfide bond formation. Further structural and chemical studies will be needed to fully characterize the inhibitory species.

\section{Materials and Methods}

Yeast strains and primers are listed in Table S2. YAP1 C-terminal tagging with EGFP and gene deletion strains were constructed as described \citep{Longtine1998,Guldener1996}. Variomics libraries were gifts from Dr. Xuewen Pan \citep{Pan2004}. For manual screening, $1$--$3 \times 10^7$ cells were plated on SC + 10 $\mu$g/mL thiolutin plates. High-throughput screens used barcode sequencing (Bar-seq). \textit{In vitro} transcription assays were performed with purified Pol II using ssDNA templates and NTPs including $^{32}$P-labeled UTP. ChIP-seq used FLAG-tagged Rpb3 with \textit{S. pombe} spike-in normalization.

\bibliography{references}

\end{document}
